A binary star system consists of $M_1$ and $M_2$ separated by a distance $D$.
$M_1$ and $M_2$ are revolving with an angular velocity $\omega$ in circular
orbits about their common centre of mass. Mass is continuously being
transferred from one star to the other. This transfer of mass causes their
orbital period and their separation to change slowly with time.

In order to simplify the analysis, we will assume that the stars are like
point particles and that the effects of the rotation about their own axes are
negligible.

\begin{description}
\item[(a)] What is the total angular momentum and kinetic energy of the
  system? \hfill(2 points)
\item[(b)] Find the relation between the angular velocity $\omega$ and the
  distance $D$ between the stars. \hfill(2 points)
\item[(c)] In a time duration $\Delta t$, a mass transfer between the two
  stars results in a change of mass $\Delta M_1$ in star $M_1$; find the
  quantity $\Delta\omega$ in terms of $\omega$, $M_1$, $M_2$ and
  $\Delta M_1$. \hfill(3 points)
\item[(d)] In a certain binary system, $M_1 = 2.9\,M_\odot$,
  $M_2 = 1.4\,M_\odot$ and the orbital period, $T = 2.49$ days. After 100
  years, the period $T$ has increased by 20\,s. Find the value of
  $\dfrac{\Delta M_1}{M_1 \Delta t}$ (in the unit ``per year'').
  \hfill(1.5 points)
\item[(e)] In which direction is mass flowing, from $M_1$ to $M_2$, or $M_2$
  to $M_1$? \hfill(0.5 point)
\item[(f)] Find also the value of $\dfrac{\Delta D}{D \Delta t}$ (in the
  unit ``per year''). \hfill(1 point)
\end{description}

You may use these approximations:
\[ (1+x)^n \approx 1 + nx, \quad \mbox{when } x \ll 1; \]
\[ (1+x)(1+y) \approx 1 + x + y, \quad \mbox{when } x \ll 1,\ y \ll 1. \]
